\chapter{Introduction}\label{ch:intro}

The central result is the Fundamental Theorem for translation-invariant
matrix product vectors. The basic object is a family of matrices
$A=\{A^i\}_{i=0}^{d-1}$, which assigns to each chain length~$N$ the vector
$\ket{V^{(N)}(A)}$ with coefficients indexed by
$\sigma=(i_1,\ldots,i_N)$,
\begin{align}
    V^{(N)}(A)_\sigma
     & = \tr(A^{i_1}\cdots A^{i_N}).
    \label{eq:intro_mpv_coefficient}
\end{align}
The main argument, for translation-invariant tensors, follows the canonical-form and
basis-of-normal-tensors argument of~\cite{Cirac2016MPDO_arXiv,Cirac2021Matrix},
with the single-block injective case anchored in the earlier representation
theorem of~\cite{PerezGarcia2007Matrix}.

\section*{Notation}

\begin{center}
    \begin{tabular}{ll}
        $\C$, $\N$, $\Z$            & complex numbers, natural numbers, integers       \\
        $\MN{D}$, $\GL_D(\C)$       & complex matrices and invertible complex matrices \\
        $\tr(X)$, $\spn_{\C}(S)$    & trace and complex linear span                    \\
        $A=\{A^i\}_{i=0}^{d-1}$     & a translation-invariant MPS tensor               \\
        $w=(i_1,\ldots,i_L)$, $A^w$ & a word and the product
        $A^{i_1}\cdots A^{i_L}$                                                        \\
        $\ket{V^{(N)}(A)}$          & the length-$N$ matrix product vector             \\
        $V^{(N)}(A)_\sigma$         & the coefficient of $\ket{V^{(N)}(A)}$ at
        $\sigma$                                                                       \\
        $O_{AB}(N)$                 & the bilinear overlap
        $\sum_\sigma V^{(N)}(A)_\sigma\,\overline{V^{(N)}(B)_\sigma}$                  \\
        $\E_A(X)$                   & the transfer map $\sum_i A^i X(A^i)^\dagger$     \\
        $A^{[L]}$                   & the length-$L$ blocked tensor                    \\
    \end{tabular}
\end{center}

In canonical form one writes
\begin{align}
    A^i
     & = \bigoplus_{j=1}^{g} \mu_j A_j^i,
    \label{eq:intro_canonical_form_a}     \\
    B^i
     & = \bigoplus_{k=1}^{h} \nu_k B_k^i.
    \label{eq:intro_canonical_form_b}
\end{align}
The blocks in (\ref{eq:intro_canonical_form_a}) and (\ref{eq:intro_canonical_form_b})
are normal tensors. If $\ket{V^{(N)}(A)}=\ket{V^{(N)}(B)}$ for every positive chain
length~$N$, then the bases of normal tensors have the same size. After relabelling the
blocks of $B$, each matched pair satisfies
\begin{align}
    B_j^i
     & = \zeta_j X_j A_j^i X_j^{-1},
    \label{eq:intro_matched_pair_gauge}
\end{align}
for some invertible matrix $X_j$ and nonzero scalar $\zeta_j$, and the equal
case compares the coefficient power sums attached to the relabelled bases. The
block gauges in (\ref{eq:intro_matched_pair_gauge}) and the weight
comparison then combine into a single gauge conjugacy of the two weighted
block-diagonal tensors in (\ref{eq:intro_canonical_form_a})
and (\ref{eq:intro_canonical_form_b}). For a single injective block this
reduces to the conclusion $B^i = X A^i X^{-1}$ for all physical
indices $i$, the single-block specialization of
(\ref{eq:intro_matched_pair_gauge}).

Chapters~\ref{ch:mps}--\ref{ch:ft_proof} carry out this argument.
Chapter~\ref{ch:mps} fixes the matrix product vector notation, word
evaluations, transfer maps, gauge equivalence, injectivity, blocking, and
site-periodic tensor families. Chapter~\ref{ch:single} proves the
single-block theorem by trace pairings, linear extension, simplicity of matrix
algebras, and Skolem--Noether. Chapters~\ref{ch:qpf}
and~\ref{ch:spectral} develop the Perron--Frobenius and
overlap-decay tools used to control normal blocks. The general theory of
Schwarz inequalities and multiplicative domains is developed in the companion
quantum-channel volume~\cite[Chapter ``Schwarz Inequalities and Multiplicative
    Domains'']{QICLean2026Blueprint}. Chapters~\ref{ch:wielandt} and~\ref{ch:canonical}
supply the blocking, primitivity, irreducibility, and canonical-form reductions.
Chapter~\ref{ch:canonical} includes the invariant-projection splitting.
Chapter~\ref{ch:bnt} supplies bases of normal tensors and the
power-sum comparison.

Chapter~\ref{ch:symmetry} uses the Fundamental Theorem to obtain virtual
symmetry gauges, projective bond representations, and string-order criteria;
its Kraus-mixing step uses the rectangular Kraus-freedom theorem from the
companion quantum-channel volume~\cite[``Rectangular Kraus
    freedom'']{QICLean2026Blueprint}.

Beyond Chapter~\ref{ch:symmetry}, subsequent chapters develop
parent Hamiltonians (Chapter~\ref{ch:parent_hamiltonian}),
correlation decay (Chapter~\ref{ch:correlations}),
concrete examples (Chapter~\ref{ch:examples}),
matrix product density operators (Chapter~\ref{ch:mpdo}),
renormalization fixed points (Chapters~\ref{ch:mps_rfp} and~\ref{ch:mpdo_rfp}),
matrix product unitaries (Chapter~\ref{ch:mpu}),
periodic tensors (Chapter~\ref{ch:periodic_ft_apps}),
the algebraic Fundamental Theorem (Chapter~\ref{ch:algebraic_ft}),
two-dimensional tensor networks (PEPS, Chapter~\ref{ch:peps_ft}),
and the asymmetric Fundamental Theorem (Chapter~\ref{ch:asymmetric_ft}).
General quantum-channel representations, positive-map examples, dynamical
semigroups, operator convexity, and entropy are developed in the companion
quantum-channel volume~\cite[Chapters ``Quantum Channels and Positive Maps,'' ``Channel Representations and Normal Forms,'' ``Positive but Not
    Completely Positive Maps,'' ``Operator Convexity and Jensen Inequalities,''
    ``Quantum Dynamical Semigroups,'' and ``Quantum Entropy'']{QICLean2026Blueprint}.
